\documentclass[11pt,letterpaper]{article}

\usepackage[utf8]{inputenc}
\usepackage[T1]{fontenc}
\usepackage[spanish,es-noshorthands]{babel}
\usepackage{spectral}
\usepackage[margin=2.5cm,headheight=48pt,headsep=18pt]{geometry}
\usepackage{fancyhdr}
\usepackage{graphicx}
\usepackage{xcolor}
\usepackage{listings}
\usepackage{tikz}
\usepackage{pgfplots}
\usepackage{booktabs}
\usepackage{enumitem}
\usepackage{float}
\usepackage[hidelinks]{hyperref}
\pgfplotsset{compat=1.18}

\newcommand{\curso}{CC4303-1 Redes}
\newcommand{\tituloinforme}{Actividad: Control de congestión}

\renewcommand{\sectionmark}[1]{\markboth{#1}{}}
\setlength{\parskip}{0.6em}
\setlength{\parindent}{0pt}

% Cabecera y pie de página con el mismo estilo del informe de práctica.
\pagestyle{fancy}
\fancyhf{}
\fancyhead[L]{\small\leftmark}
\fancyfoot[L]{\small Informe \tituloinforme}
\fancyfoot[R]{\small CC4303-1}
\renewcommand{\headrulewidth}{0.4pt}
\renewcommand{\footrulewidth}{0.4pt}

\fancypagestyle{portada}{%
  \fancyhf{}
  \fancyhead[L]{\small\begin{tabular}[b]{@{}l@{}}
    Universidad de Chile\\
    Facultad de Ciencias Físicas y Matemáticas\\
    Departamento de Ciencias de la Computación\\
    \curso
  \end{tabular}}
  \fancyhead[R]{\includegraphics[height=1.3cm]{logo_fcfm.png}}
  \renewcommand{\headrulewidth}{0.4pt}
  \renewcommand{\footrulewidth}{0pt}
}

\definecolor{fondocodigo}{gray}{0.96}
\lstset{
  basicstyle=\ttfamily\footnotesize,
  backgroundcolor=\color{fondocodigo},
  frame=single,
  rulecolor=\color{gray!50},
  breaklines=true,
  columns=fullflexible,
  keepspaces=true,
  showstringspaces=false,
  literate={á}{{\'a}}1 {é}{{\'e}}1 {í}{{\'i}}1 {ó}{{\'o}}1 {ú}{{\'u}}1
           {Á}{{\'A}}1 {ñ}{{\~n}}1
}

\begin{document}

\begin{titlepage}
  \thispagestyle{portada}
  \centering
  \vspace*{5cm}
  {\LARGE Informe\par}
  \vspace{0.3cm}
  {\Large \tituloinforme\par}
  \vspace{0.4cm}
  {\large Juan Pablo Bustamante Flaño\par}
  \vfill
  Correo electrónico: juan.bustamante.f@ug.uchile.cl\\
  Curso: \curso{} --- Primavera 2026\\
  Profesora: Ivana Bachmann\\
  Módulo 3: Protocolos de transporte, semanas 8-10\\
  Fecha: 19/10/2026
  \vspace*{3cm}
\end{titlepage}

\section{Introducción}

En esta actividad se extendió el TCP simplificado de la actividad de Stop \& Wait con
Go-Back-N y con control de congestión basado en TCP Tahoe (solo slow start y congestion
avoidance). El trabajo se hizo sobre una copia del código anterior.

\begin{center}
\begin{tabular}{@{}ll@{}}
  \toprule
  Archivo & Contenido \\
  \midrule
  \texttt{CongestionControl.py} & Clase \texttt{CongestionControl} (parte 1). \\
  \texttt{SocketTCP.py} & \texttt{SocketTCP} con Stop \& Wait y Go-Back-N (partes 2 y 3). \\
  \texttt{socketUDP.py}, \texttt{slidingWindowCC.py} & Sin modificar. \\
  \texttt{CongestionControl\_test.py} & Sin modificar. \\
  \texttt{cliente.py}, \texttt{servidor.py} & Cliente y servidor de prueba. \\
  \texttt{test\_cliente.py}, \texttt{test\_servidor.py} & Pruebas de \texttt{send}/\texttt{recv} de la actividad anterior. \\
  \texttt{experimento.sh} & Comparación con y sin control de congestión. \\
  \bottomrule
\end{tabular}
\end{center}

Uso:

\begin{lstlisting}
% python3 servidor.py localhost 8000 [stop_and_wait|go_back_n] > recibido.txt
% python3 cliente.py localhost 8000 [stop_and_wait|go_back_n] [--sin-cc] < archivo.txt
\end{lstlisting}

El modo por defecto es \texttt{go\_back\_n} con control de congestión activado. La
opción \texttt{-{}-sin-cc} lo desactiva. El cliente informa en \texttt{stderr} el tiempo
que tomó \texttt{send} y la cantidad de segmentos enviados. El timeout
(\texttt{TIMEOUT = 1.5}) y el modo debug (\texttt{DEBUG = True}) son variables globales
al comienzo de \texttt{SocketTCP.py}. El header de los segmentos es el mismo de la actividad anterior:
texto con el formato \texttt{SYN|ACK|FIN|SEQ|} y \texttt{SEQ} de 10 dígitos (17 bytes).

\section{Parte 1: clase \texttt{CongestionControl}}

El constructor \texttt{CongestionControl(MSS)} inicializa \texttt{current\_state} en slow
start, \texttt{cwnd} en 1 MSS y \texttt{ssthresh} en \texttt{None}. Las funciones son:

\begin{itemize}[nosep]
  \item \texttt{get\_cwnd()}: tamaño de la ventana en bytes.
  \item \texttt{get\_MSS\_in\_cwnd()}: cantidad de MSS completos en \texttt{cwnd}
        ($\lfloor \texttt{cwnd}/\texttt{MSS} \rfloor$; con 17 bytes y MSS = 4 retorna 4).
  \item \texttt{event\_ack\_received()}: en slow start suma 1 MSS y, si
        \texttt{ssthresh} ya está definido y \texttt{cwnd >= ssthresh}, pasa a congestion
        avoidance. En congestion avoidance suma
        $\texttt{MSS}/\texttt{get\_MSS\_in\_cwnd()}$, es decir, 1 MSS por cada ventana
        completa de ACKs (Additive Increase).
  \item \texttt{event\_timeout()}: en ambos estados fija
        \texttt{ssthresh = int(cwnd/2)}, deja \texttt{cwnd} en 1 MSS y vuelve a slow start
        (Multiplicative Decrease).
  \item \texttt{is\_state\_slow\_start()}, \texttt{is\_state\_congestion\_avoidance()} y
        \texttt{get\_ssthresh()}.
\end{itemize}

\textbf{Decisiones de diseño.}
\begin{itemize}[nosep]
  \item En congestion avoidance \texttt{cwnd} suma fracciones de MSS (por ejemplo,
        $8/3$), así que deja de ser entero y se guarda como \texttt{float}. Con números de
        punto flotante, la suma de tres incrementos de $8/3$ puede quedar en
        $31{,}999\ldots$ en vez de 32, y \texttt{get\_MSS\_in\_cwnd()} entregaría un MSS
        menos. Para evitarlo, el cociente se redondea a 6 decimales antes de truncarlo:
        \texttt{int(round(cwnd / MSS, 6))}.
  \item \texttt{ssthresh} se trunca a entero (con el mismo redondeo previo), que es lo que
        espera el test de u-cursos (\texttt{int(5/2*MSS)}).
  \item El cambio de estado se revisa solo al aumentar \texttt{cwnd}, como indica el
        paso 4. Por eso, después de un timeout con \texttt{cwnd} = 1 MSS, el
        \texttt{ssthresh} queda en $\texttt{MSS}/2$ y el primer ACK lleva directamente a
        congestion avoidance.
\end{itemize}

\textbf{Test.} \texttt{CongestionControl\_test.py} pasó con MSS = 7 (el valor del
archivo) y también con 1, 3, 8 y 13:

\begin{lstlisting}
% python3 CongestionControl_test.py
Initial test passed
ACK test passed
Timeout test passed
ssthresh test passed
Timeout test 2 passed
\end{lstlisting}

\section{Parte 2: Go-Back-N}

\subsection{Preguntas previas}

\textbf{(1) ¿Por qué sirve el código de Stop \& Wait como base para Go-Back-N? ¿Qué
similitud tienen?} Stop \& Wait es un caso particular de Go-Back-N con ventana de tamaño
1. Ambos usan un solo timer, ACKs acumulativos (el ACK indica el siguiente byte esperado)
y la misma reacción ante un timeout: retransmitir todo lo enviado que aún no está
confirmado. El código de ejemplo ya usa \texttt{SlidingWindowCC} con ventana 1 y el timer
de \texttt{SocketUDP}. Para pasar a Go-Back-N basta con agrandar la ventana, enviar todos
sus elementos, avanzarla según el ACK acumulado y reenviar la ventana completa cuando
vence el timer.

\textbf{(2) ¿Qué pasa con \texttt{recv}? ¿Cambia al usar \texttt{SocketUDP} y
\texttt{SlidingWindowCC}?} En Go-Back-N el receptor no cambia respecto de Stop \& Wait.
Acepta solo el segmento cuyo \texttt{seq} es el esperado, descarta los repetidos o fuera de
orden y siempre responde con el ACK acumulado. No guarda segmentos fuera de orden (a
diferencia de Selective Repeat), así que no necesita una ventana. Tampoco necesita
timers, porque el receptor nunca retransmite. Por eso \texttt{recv\_using\_go\_back\_n}
llama a la misma lógica que \texttt{recv\_using\_stop\_and\_wait}. Mi receptor de Stop \&
Wait ya reenviaba el ACK acumulado ante cualquier segmento de datos inesperado, así que
sirve sin cambios. Lo único que cambia es el socket: ahora se recibe con
\texttt{SocketUDP}.

\subsection{Uso de \texttt{SocketUDP} y \texttt{SlidingWindowCC}}

Ambas clases se importan tal como vienen de u-cursos y se usan solo a través de sus
funciones documentadas (\texttt{settimeout}, \texttt{set\_timer\_list\_length},
\texttt{sendto}, \texttt{recvfrom}, \texttt{stop\_timer} y \texttt{close}; y en la
ventana, \texttt{get\_data}, \texttt{get\_sequence\_number}, \texttt{move\_window} y
\texttt{update\_window\_size}). Toda la comunicación de \texttt{SocketTCP} (handshake,
Stop \& Wait, Go-Back-N y cierre) pasa ahora por \texttt{SocketUDP}, siguiendo el ejemplo
de Stop \& Wait de la sección \emph{SocketUDP y SlidingWindowCC}: se envía con
\texttt{sendto}, se espera con \texttt{recvfrom} dentro de un \texttt{try} y, al recibir
el ACK válido, se detiene el timer con \texttt{stop\_timer}. Usar esta clase obligó a
tomar dos decisiones:

\begin{itemize}[nosep]
  \item \textbf{Un segundo timer para los ACK.} \texttt{sendto} inicia el timer indicado
        en cada envío. Si un ACK usara el timer 0, más adelante lanzaría un
        \texttt{TimeoutError} falso, aunque un ACK no espera respuesta. Por eso se usan
        dos timers (\texttt{set\_timer\_list\_length(2)}): el timer 0 para los segmentos
        que esperan respuesta (datos, \texttt{SYN}, \texttt{SYN+ACK}, \texttt{FIN} y
        \texttt{FIN+ACK}) y el timer 1 para los ACK, que se detiene con
        \texttt{stop\_timer(1)} justo después de enviarlos.
  \item \textbf{Detener el timer antes de reiniciarlo.} \texttt{SocketUDP} usa
        \texttt{threading.Timer}, que no se puede iniciar dos veces. Si el timer venció y
        se vuelve a llamar a \texttt{sendto} sin \texttt{stop\_timer}, se lanza
        \texttt{RuntimeError}. Por eso, después de cada \texttt{TimeoutError} y antes de
        cada retransmisión, se llama a \texttt{stop\_timer(0)}. Además, \texttt{stop\_timer}
        falla si el timer no existe, así que \texttt{SocketTCP} guarda en
        \texttt{data\_timer\_on} si el timer 0 está activo y solo lo detiene en ese caso.
        Por la misma razón, mientras se espera un ACK no se reenvía nada ante segmentos
        inesperados: solo se retransmite después de un timeout.
\end{itemize}

\subsection{Funciones \texttt{send} y \texttt{recv}}

El código anterior se movió a \texttt{send\_using\_stop\_and\_wait} y
\texttt{recv\_using\_stop\_and\_wait}. \texttt{send(message, mode)} y
\texttt{recv(buff\_size, mode)} eligen la implementación con \texttt{mode}, que acepta
\texttt{"stop\_and\_wait"} (por defecto) o \texttt{"go\_back\_n"}.

\texttt{send\_using\_go\_back\_n(message)} funciona así:

\begin{enumerate}[nosep]
  \item Arma \texttt{data\_list} con el largo del mensaje seguido del mensaje en trozos de
        MSS = 8 bytes, y crea la ventana \texttt{SlidingWindowCC(window\_size, data\_list,
        self.seq)}.
  \item Envía todos los elementos de la ventana. El timer parte con el primero.
  \item Espera un ACK válido, es decir, con \texttt{seq} mayor que el del primer elemento
        de la ventana y menor o igual que el final del mensaje. Los ACK repetidos se
        ignoran.
  \item Con un ACK válido, avanza la ventana en la cantidad de elementos que confirma y
        envía seguidos los elementos que entraron a ella. Si hubo envíos nuevos, el timer
        se reinicia; si no, sigue corriendo para cubrir los segmentos que todavía no se
        confirman.
  \item Si vence el timer, reenvía la ventana completa.
\end{enumerate}

Para saber qué elementos de la ventana todavía no se envían, se guarda el \texttt{seq}
siguiente al último segmento enviado.

\textbf{Test sin pérdida.} Las pruebas de la actividad anterior (con el test 4 de
$2n = 42$ bytes y \texttt{recv(21)}) pasaron con \texttt{mode="go\_back\_n"}. El envío de
archivos de 304 y 1201 bytes llegó idéntico según \texttt{cmp}.

\textbf{Test con pérdida.} Se repitió induciendo pérdidas con \texttt{netem} (ver
sección~\ref{sec:netem}). Los resultados están en la sección~\ref{sec:pruebas}.

\section{Parte 3: control de congestión}

\begin{itemize}[nosep]
  \item \texttt{congestion\_controler} es un \texttt{CongestionControl(8)} que se crea en
        el constructor de \texttt{SocketTCP}. Se mantiene entre llamadas a
        \texttt{send}, igual que la \texttt{cwnd} de TCP se mantiene durante la conexión.
  \item \texttt{congestion\_control\_on} se activa o desactiva con el setter
        \texttt{set\_congestion\_control(bool)}. Por defecto vale \texttt{True}.
  \item Con control de congestión, la ventana parte con
        \texttt{window\_size = get\_MSS\_in\_cwnd()}. Sin él, la ventana se mantiene fija
        en 8 segmentos. En ambos casos el mensaje se divide en trozos de 8 bytes, para que
        la comparación de la sección~\ref{sec:experimento} sea justa.
  \item Cada ACK válido llama a \texttt{event\_ack\_received()} y cada timeout llama a
        \texttt{event\_timeout()}. Después de cada evento se recalcula
        \texttt{window\_size} y se aplica con \texttt{data\_window.update\_window\_size}.
        Si la ventana creció, sus nuevos elementos se envían seguidos.
\end{itemize}

\textbf{Valor inicial de \texttt{window\_size}.} Debe ser 1, porque \texttt{cwnd} parte en
1 MSS y \texttt{get\_MSS\_in\_cwnd()} retorna 1. El modo debug lo confirma:

\begin{lstlisting}
(debug) [CC] inicio: estado=slow start, cwnd=8 B, ssthresh=None, window_size=1
\end{lstlisting}

\textbf{ACK que se cuenta una vez.} \texttt{event\_ack\_received()} se llama una vez por
cada ACK que avanza la ventana, aunque ese ACK confirme varios segmentos porque se
perdieron ACKs anteriores. Así se cuentan ACKs y no bytes confirmados, como en TCP Tahoe
sin \emph{byte counting}. Los ACK repetidos no generan eventos, porque esta actividad no
implementa fast retransmit.

\textbf{Caso borde: la ventana se achica y llega un ACK fuera de ella.} Después de un
timeout la ventana queda en 1 elemento, pero puede llegar un ACK de segmentos enviados con
la ventana anterior, cuyo \texttt{seq} supera el último elemento de la ventana actual. En
ese caso se avanza la ventana completa, y se repite hasta que el ACK caiga dentro de ella.
Luego se envían los elementos que no estaban confirmados. El caso ocurrió 5 veces en las
pruebas con \texttt{netem}. En el siguiente ejemplo, la ventana de 2 elementos tenía los
segmentos \texttt{seq=835} y \texttt{843}. El receptor recibió ambos, pero sus ACK no
llegaron al emisor. Tras el timeout la ventana quedó solo con \texttt{seq=835}, y el ACK
acumulado \texttt{seq=851}, que llegó en el tercer intento, confirma también el segmento
\texttt{seq=843}, que ya no estaba en la ventana:

\begin{lstlisting}
# emisor
(debug) [CC] ACK: estado=congestion avoidance, cwnd=16 B, ssthresh=8, window_size=2
(debug) Enviando [DATA seq=835 data=b'sum dolo']
(debug) Enviando [DATA seq=843 data=b'r sit am']
(debug) Timeout esperando ACK de seq=835, retransmitiendo la ventana
(debug) [CC] timeout: estado=slow start, cwnd=8 B, ssthresh=8, window_size=1
(debug) Enviando [DATA seq=835 data=b'sum dolo']
(debug) Timeout esperando ACK de seq=835, retransmitiendo la ventana
(debug) [CC] timeout: estado=slow start, cwnd=8 B, ssthresh=4, window_size=1
(debug) Enviando [DATA seq=835 data=b'sum dolo']
(debug) Timeout esperando ACK de seq=835, retransmitiendo la ventana
(debug) [CC] timeout: estado=slow start, cwnd=8 B, ssthresh=4, window_size=1
(debug) Enviando [DATA seq=835 data=b'sum dolo']
(debug) Recibido [ACK seq=851]
(debug) ACK seq=851 más allá de la ventana, se avanzó 2 veces
(debug) [CC] ACK: estado=congestion avoidance, cwnd=16 B, ssthresh=4, window_size=2
(debug) Enviando [DATA seq=851 data=b'et, cons']
(debug) Enviando [DATA seq=859 data=b'ectetur ']
# receptor
(debug) Recibido [DATA seq=835 data=b'sum dolo'], enviando ACK seq=843
(debug) Recibido [DATA seq=843 data=b'r sit am'], enviando ACK seq=851
(debug) Recibido [DATA seq=835 data=b'sum dolo'] fuera de orden, reenviando ACK seq=851
\end{lstlisting}

\section{Pruebas}
\label{sec:pruebas}

\subsection{Pérdidas con \texttt{netem}}
\label{sec:netem}

Las pruebas con pérdida se hicieron con \texttt{netem}, con la configuración del
enunciado:

\begin{lstlisting}
% sudo tc qdisc add dev lo root netem loss 20.0% delay 0.5s
\end{lstlisting}

Con este delay el RTT es de 1 s, menor que el timeout de 1,5 s que se usó en estas
pruebas.

\subsection{Integridad de los datos}

\begin{center}
\begin{tabular}{@{}lcc@{}}
  \toprule
  Prueba (\texttt{netem}: 20\% de pérdida, delay 0,5 s) & Criterio & Resultado \\
  \midrule
  Archivo de 1201 B, Go-Back-N con control de congestión & \texttt{cmp} idéntico & 3/3 \\
  Archivo de 1201 B, Go-Back-N sin control de congestión & \texttt{cmp} idéntico & 3/3 \\
  Archivo de 304 B, Stop \& Wait & \texttt{cmp} idéntico & 2/2 \\
  Tests 1--4 de \texttt{send}/\texttt{recv}, Go-Back-N & 4 de 4 \texttt{Passed} & 3/3 \\
  \bottomrule
\end{tabular}
\end{center}

Con el delay real, enviar el archivo de 1201 B tomó en promedio 153 s con control de
congestión (138, 178 y 145 s; 474 segmentos intercambiados en promedio) y 81 s sin él
(82, 75 y 86 s; 723 segmentos). Es la misma tendencia del experimento de la
sección~\ref{sec:experimento}. El caso borde de la ventana que se achica (parte 3)
ocurrió 5 veces en estas ejecuciones.

\subsection{Comportamiento sin pérdida}

Sin pérdida inducida no ocurrieron timeouts espontáneos: en \texttt{localhost} el RTT es
del orden de microsegundos, muy inferior a cualquier timeout usado. Por eso el control de
congestión se queda siempre en slow start: \texttt{ssthresh} sigue en \texttt{None} y la
ventana crece en 1 MSS por ACK. Con 1201 bytes (152 segmentos de datos), \texttt{cwnd}
llega a $153 \cdot 8$ bytes y el envío toma unos 8 ms:

\begin{lstlisting}
(debug) [CC] inicio: estado=slow start, cwnd=8 B, ssthresh=None, window_size=1
(debug) [CC] ACK: estado=slow start, cwnd=16 B, ssthresh=None, window_size=2
(debug) [CC] ACK: estado=slow start, cwnd=24 B, ssthresh=None, window_size=3
(debug) [CC] ACK: estado=slow start, cwnd=32 B, ssthresh=None, window_size=4
...
STATS cliente: tiempo_envio=0.008s segmentos_enviados=154
\end{lstlisting}

Esto es lo esperado: sin pérdidas no hay eventos que reduzcan la ventana ni que definan
\texttt{ssthresh}.

\subsection{Comportamiento con pérdida}

Con \texttt{netem} (20\% de pérdida y delay de 0,5 s), cada timeout define
\texttt{ssthresh} como la mitad de \texttt{cwnd} y devuelve la ventana a 1 MSS. Al
alcanzar \texttt{ssthresh} se pasa a congestion avoidance, donde la ventana crece en
$8/n$ bytes por ACK. Este extracto es de una de las ejecuciones con el archivo de 1201 B:

\begin{lstlisting}
(debug) [CC] timeout: estado=slow start, cwnd=8 B, ssthresh=8, window_size=1
(debug) [CC] ACK: estado=congestion avoidance, cwnd=16 B, ssthresh=8, window_size=2
(debug) [CC] ACK: estado=congestion avoidance, cwnd=20 B, ssthresh=8, window_size=2
(debug) [CC] ACK: estado=congestion avoidance, cwnd=24 B, ssthresh=8, window_size=3
(debug) [CC] ACK: estado=congestion avoidance, cwnd=26.6667 B, ssthresh=8, window_size=3
(debug) [CC] ACK: estado=congestion avoidance, cwnd=29.3333 B, ssthresh=8, window_size=3
(debug) [CC] ACK: estado=congestion avoidance, cwnd=32 B, ssthresh=8, window_size=4
...
(debug) [CC] ACK: estado=congestion avoidance, cwnd=40 B, ssthresh=8, window_size=5
(debug) [CC] ACK: estado=congestion avoidance, cwnd=41.6 B, ssthresh=8, window_size=5
...
(debug) [CC] ACK: estado=congestion avoidance, cwnd=48 B, ssthresh=8, window_size=6
(debug) Timeout esperando ACK de seq=243, retransmitiendo la ventana
(debug) [CC] timeout: estado=slow start, cwnd=8 B, ssthresh=24, window_size=1
(debug) [CC] ACK: estado=slow start, cwnd=16 B, ssthresh=24, window_size=2
\end{lstlisting}

Los cambios son correctos. Con \texttt{cwnd} = 8 B, el timeout deja
\texttt{ssthresh} $= 8/2 = 4$ o, como en la primera línea, 8 cuando \texttt{cwnd} era
16 B. El siguiente ACK supera \texttt{ssthresh} y pasa a congestion avoidance, donde con
16 B se suman $8/2 = 4$ bytes por ACK, con 24 B $8/3$, con 32 B $8/4 = 2$ y con 40 B
$8/5 = 1{,}6$. El timeout con \texttt{cwnd} = 48 B deja \texttt{ssthresh} = 24 y la
ventana vuelve a slow start, donde se suman 8 bytes por ACK. La figura~\ref{fig:cwnd}
muestra los primeros 60 eventos de esa ejecución.

\begin{figure}[H]
\centering
\begin{tikzpicture}
\begin{axis}[
  width=\textwidth, height=6cm,
  xlabel={Evento de control de congestión}, ylabel={\texttt{cwnd} [bytes]},
  xmin=0, xmax=59, ymin=0, ymax=52,
  ytick={0,8,16,24,32,40,48},
  grid=major, grid style={gray!20},
  legend style={at={(0.5,1.03)},anchor=south,legend columns=3,font=\footnotesize},
]
\addplot[thick,blue!60!black,mark=*,mark size=1pt] coordinates {
(0,8) (1,8) (2,8) (3,8) (4,8) (5,16) (6,8) (7,8) (8,16) (9,8) (10,16) (11,20) (12,24)
(13,26.67) (14,29.33) (15,32) (16,34) (17,36) (18,38) (19,40) (20,41.6) (21,43.2)
(22,44.8) (23,46.4) (24,48) (25,8) (26,16) (27,8) (28,16) (29,20) (30,24) (31,8) (32,16)
(33,20) (34,24) (35,26.67) (36,29.33) (37,32) (38,34) (39,36) (40,38) (41,40) (42,41.6)
(43,43.2) (44,44.8) (45,46.4) (46,48) (47,8) (48,16) (49,24) (50,26.67) (51,29.33)
(52,32) (53,34) (54,36) (55,38) (56,40) (57,41.6) (58,8) (59,16)};
\addlegendentry{\texttt{cwnd}}
\addplot[only marks,red!70!black,mark=x,mark size=3pt,thick] coordinates {
(1,8) (2,8) (3,8) (4,8) (6,8) (7,8) (9,8) (25,8) (27,8) (31,8) (47,8) (58,8)};
\addlegendentry{timeout}
\addplot[const plot,dashed,orange!80!black,thick] coordinates {
(1,4) (6,8) (7,4) (9,8) (25,24) (27,8) (31,12) (47,24) (58,20) (59,20)};
\addlegendentry{\texttt{ssthresh}}
\end{axis}
\end{tikzpicture}
\caption{Evolución de \texttt{cwnd} y \texttt{ssthresh} con \texttt{netem} (20\% de
pérdida, delay 0,5 s, MSS = 8).}
\label{fig:cwnd}
\end{figure}

Se observa el patrón característico de Tahoe: crecimiento rápido en slow start,
crecimiento lineal más lento en congestion avoidance y caída a 1 MSS en cada timeout. Con
20\% de pérdida en cada dirección, la ventana nunca superó los 6 MSS (48 B) en esta
ejecución.

\subsection{Con versus sin control de congestión}
\label{sec:experimento}

El enunciado pide un timeout menor a 1 s (0,5 s). Con el delay de 0,5 s de
\texttt{netem}, el RTT de 1 s superaría ese timeout y todos los segmentos se
retransmitirían antes de recibir su ACK. Por eso, para este experimento \texttt{netem} se
configuró solo con pérdida:

\begin{lstlisting}
% sudo tc qdisc change dev lo root netem loss 20.0%
\end{lstlisting}

\texttt{experimento.sh} transmite un archivo 5 veces con Go-Back-N con y sin control de
congestión. Antes de ejecutarlo se configura \texttt{TIMEOUT = 0.5} y
\texttt{DEBUG = False} en \texttt{SocketTCP.py}:

\begin{lstlisting}
% bash experimento.sh archivo_100.txt 5
% bash experimento.sh archivo_1200.txt 5
\end{lstlisting}

Para cada ejecución registra el tiempo de \texttt{send} y la cantidad de segmentos
enviados por el cliente (datos, más \texttt{SYN} y \texttt{ACK} del handshake) y por el
servidor (\texttt{SYN+ACK} y ACKs). Siguiendo el enunciado, se probó un archivo de 100 B
y, para responder la pregunta sobre archivos de más de 1000 B, uno de 1201 B. En todas
las ejecuciones el archivo llegó íntegro.

\begin{center}
\begin{tabular}{@{}lcccc@{}}
  \toprule
  & \multicolumn{2}{c}{100 B} & \multicolumn{2}{c}{1201 B} \\
  \cmidrule(lr){2-3}\cmidrule(lr){4-5}
  Rep. & con CC & sin CC & con CC & sin CC \\
  \midrule
  1 & 1,502 s / 43 & 3,002 s / 98 & 29,529 s / 487 & 17,521 s / 765 \\
  2 & 2,503 s / 59 & 1,501 s / 50 & 19,525 s / 465 & 18,022 s / 715 \\
  3 & 2,503 s / 42 & 2,002 s / 60 & 20,522 s / 458 & 24,025 s / 892 \\
  4 & 3,003 s / 57 & 1,501 s / 52 & 29,530 s / 506 & 21,023 s / 862 \\
  5 & 2,503 s / 42 & 3,002 s / 100 & 31,528 s / 491 & 15,020 s / 674 \\
  \midrule
  Promedio & \textbf{2,403 s / 48,6} & \textbf{2,202 s / 72,0} & \textbf{26,127 s / 481,4} & \textbf{19,122 s / 781,6} \\
  \bottomrule
\end{tabular}

\smallskip
{\footnotesize Cada celda: tiempo de \texttt{send} / segmentos totales intercambiados
(cliente + servidor).}
\end{center}

\textbf{¿Cuál toma menos tiempo?} En estas condiciones, Go-Back-N \emph{sin} control de
congestión. Con el archivo de 1201 B tardó en promedio 19,1 s, contra 26,1 s con control
de congestión (27\% menos). Con 100 B la diferencia es mínima (2,2 s contra 2,4 s) y queda
dentro de la variabilidad entre repeticiones. Las pruebas con delay de la
sección~\ref{sec:netem} muestran la misma tendencia (81 s contra 153 s).

\textbf{¿Cómo se comparan?} El control de congestión intercambió muchos menos segmentos:
38\% menos con 1201 B (481 contra 782) y 33\% menos con 100 B (49 contra 72). La
explicación es la siguiente:

\begin{itemize}[nosep]
  \item Sin delay, el RTT en \texttt{localhost} es casi cero, así que el tiempo total
        depende casi solo de cuántos timeouts de 0,5 s ocurren. Por eso los tiempos son
        casi múltiplos exactos de 0,5 s.
  \item Las pérdidas son aleatorias y no se deben a congestión, pero el control de
        congestión las interpreta como congestión. Después de cada timeout, la ventana
        vuelve a 1 MSS y \texttt{ssthresh} baja. Como con 20\% de pérdida casi siempre
        vuelve a ocurrir otro timeout, la ventana se mantiene pequeña (figura~\ref{fig:cwnd})
        y cada timeout permite avanzar pocos segmentos.
  \item Sin control de congestión la ventana siempre es de 8. Cada timeout retransmite
        hasta 8 segmentos, de los que muchos ya habían llegado (más tráfico), pero en cada
        intervalo entre timeouts se avanza más. El resultado son menos timeouts en total,
        aunque con más segmentos enviados.
\end{itemize}

El control de congestión no está pensado para minimizar el tiempo de un flujo aislado con
pérdidas aleatorias, sino para que la red sea justa: al reducir la ventana ante una
pérdida, el flujo deja de sobrecargar la red y le deja capacidad a otros flujos. En una
red real con congestión, enviar 62\% más segmentos (como Go-Back-N sin control) empeoraría
la congestión y aumentaría las pérdidas para todos, mientras que el control de congestión
ayuda a evitar el colapso de la red.

\section{Casos borde y limitaciones}

\begin{itemize}
  \item \textbf{Ventana que se achica.} Se resuelve avanzando la ventana completa hasta
        que el ACK caiga dentro de ella (parte 3).
  \item \textbf{Timer único de la ventana.} Cuando llega un ACK que no permite enviar
        segmentos nuevos, el timer no se reinicia. Así puede vencer antes de un timeout
        completo contado desde el último ACK, lo que solo provoca una retransmisión
        anticipada que Go-Back-N maneja correctamente.
  \item \textbf{\texttt{ssthresh} menor que 1 MSS.} Con \texttt{cwnd} = 8 B, un timeout
        deja \texttt{ssthresh} = 4 B, por lo que el siguiente ACK pasa directamente a
        congestion avoidance. Es la consecuencia de aplicar la regla
        \texttt{ssthresh = cwnd/2} sin un mínimo; TCP real usa un mínimo de 2 MSS.
  \item \textbf{Sin fast retransmit.} Los ACK repetidos se ignoran, como pide el
        enunciado (solo slow start y congestion avoidance). Implementar fast retransmit
        reduciría los timeouts y mejoraría el tiempo con control de congestión.
  \item \textbf{Cierre.} Se mantiene la limitación de la actividad anterior: si se pierden
        los 3 \texttt{FIN}, \texttt{recv\_close()} queda esperando.
  \item \textbf{Uso de CPU.} \texttt{SocketUDP.recvfrom} revisa un socket no bloqueante
        en un loop hasta que llega algo o vence un timer, así que mientras un extremo
        espera usa un núcleo de la CPU completo. Es el comportamiento de la clase
        entregada, que se usa sin modificar.
\end{itemize}

\section*{Declaración de uso de IA}

En este informe se utilizó el modelo de IA Claude exclusivamente como apoyo en la
revisión del código, recomendaciones de diseño, la redacción del informe, la creación de
las tablas y la elaboración de los gráficos y diagramas en \LaTeX{}. Las decisiones de
diseño, la implementación del código, la ejecución de los experimentos y la obtención de
los datos fueron realizadas por el autor. Todo el contenido generado con apoyo de la
herramienta fue revisado y validado, y su responsabilidad es del autor.

\end{document}
